% Section 1.4, from lectures/L01/appendix/c_exercises.md.

\section{Exercises}
\label{c1:sec:exercises}\label{c1:app:c}
These exercises reinforce the concepts from \secref{c1:app:b}.

\begin{aside}
\textbf{How to check your work.} Write each exercise set as its own application under
\file{lectures/L01/exercises}, with the file names given there, and run \code{make test} from the
repository root. A set’s tests switch on as soon as its files exist, and a set you have not started
is reported as \code{SKIP}. Each application grows through its set, and the tests check where it
ends up. Programs run on \code{native_sim} in simulated time, so their output, timestamps included,
is compared with the expected output line by line.

The solutions are published in the companion repository, under
\file{lectures/L01/appendix/solutions}, a few days after the lecture. Try each exercise yourself
before you read them.
\end{aside}

Remember \code{source env.sh} in every new terminal. Build and run an exercise from the repository
root:

\begin{shell}
west build -p -b native_sim/native/64 lectures/L01/exercises/exercise1
west build -t run
\end{shell}

% ---------------------------------------------------------------------
\exerciseset{A first application}
\label{c1:set:1}
Create a directory named \code{exercise1} to work in for this exercise set.

\exercise{The three files}{Code}
\label{c1:ex:1.1}
In this exercise you will create the smallest Zephyr application that builds.

\task{Tasks}
Create these files in \code{exercise1}:

\begin{itemize}
  \item \code{CMakeLists.txt}: require CMake 3.20, find the Zephyr package, name the project
    \code{exercise1}, and add \code{src/main.c} to the \code{app} target.
  \item \code{prj.conf}: an empty file, or one holding only a comment.
  \item \code{src/main.c}: a \code{main()} that prints \code{Hello, Zephyr!} with \code{printk()} and
    returns \code{0}.
\end{itemize}

Build and run the application.

\task{Reflection}
\begin{itemize}
  \item What happens if you delete \code{prj.conf} and build again with \code{-p}?
  \item The program prints its greeting and then keeps running. Why does it not exit when \code{main()}
    returns? How do you stop it?
  \item Where in \code{build/} is the program you ran?
\end{itemize}

\exercise{The uptime}{Code}
\label{c1:ex:1.2}

\task{Tasks}
Extend \code{main()} so that after the greeting it prints the uptime three times, with
\code{k_msleep(1000)} after each print:

\begin{console}
Hello, Zephyr!
Uptime: 0 ms
Uptime: 1010 ms
Uptime: 2020 ms
\end{console}

Use \code{k_uptime_get()}, which returns an \code{int64_t}, printed with \code{\%lld}.

Run the program with \code{./build/zephyr/zephyr.exe -stop_at=5}, and then again with
\code{-stop_at=5 -no-rt}.

\task{Reflection}
\begin{itemize}
  \item Why is the second uptime 1010 ms and not 1000 ms? Find the option in \code{build/zephyr/.config}
    that decides it.
  \item What is the difference between the two runs? Is the output different?
  \item Would the uptimes be the same if the host were busy compiling something else at the same time? Why?
\end{itemize}

% ---------------------------------------------------------------------
\exerciseset{Logging}
\label{c1:set:2}
Create an application named \code{exercise2} to work in for this exercise set, with the same three
files as \code{exercise1}.

\exercise{A temperature monitor}{Code}
\label{c1:ex:2.1}
In this exercise you will replace \code{printk()} with the logging subsystem.

\task{Tasks}
\begin{itemize}
  \item Enable logging in \code{prj.conf}.
  \item Register a log module named \code{monitor} \textbf{without} giving a level:
    \code{LOG_MODULE_REGISTER(monitor);}.
  \item In \code{main.c}, define the readings as an array:
\begin{ccode}
static const int readings[] = {21, 24, 28, 33, 36, 29};
\end{ccode}
  \item In \code{main()}, go through the readings, and for each one log:
\begin{itemize}
  \item \code{Temperature <t> C is critical} as an error, if it is 35 or more.
  \item \code{Temperature <t> C is high} as a warning, if it is 30 or more.
  \item \code{Temperature <t> C} as information, otherwise.
\end{itemize}
  \item Sleep 500 ms after each reading.
\end{itemize}

The output must be:

\begin{console}
[00:00:00.000,000] <inf> monitor: Temperature 21 C
[00:00:00.510,000] <inf> monitor: Temperature 24 C
[00:00:01.020,000] <inf> monitor: Temperature 28 C
[00:00:01.530,000] <wrn> monitor: Temperature 33 C is high
[00:00:02.040,000] <err> monitor: Temperature 36 C is critical
[00:00:02.550,000] <inf> monitor: Temperature 29 C
\end{console}

\note[Hint] \code{ARRAY_SIZE(readings)}, from \code{<zephyr/sys/util.h>}, which
\code{<zephyr/kernel.h>} includes, gives the number of elements.

\exercise{Only what matters}{Code}
\label{c1:ex:2.2}

\task{Tasks}
Without changing any file, build the application so that it logs only warnings and errors, by
setting the default log level on the command line:

\begin{shell}
west build -p -b native_sim/native/64 lectures/L01/exercises/exercise2 -- -DCONFIG_LOG_DEFAULT_LEVEL=2
\end{shell}

\noindent The output must now be only:

\begin{console}
[00:00:01.530,000] <wrn> monitor: Temperature 33 C is high
[00:00:02.040,000] <err> monitor: Temperature 36 C is critical
\end{console}

\task{Reflection}
\begin{itemize}
  \item Why did this change the monitor’s output, when the command line said nothing about \code{monitor}?
  \item Would it still work if the module were registered as \code{LOG_MODULE_REGISTER(monitor,
    LOG_LEVEL_INF)}?
  \item Are the information messages still in the program, only not printed? Explain how you know.
  \item After this build, run \code{west build -t run} a few times, and then build again \textbf{without}
    \code{-p} and without the \code{-D} option. Which log level does that build use? Look for the answer
    in \code{build/zephyr/.config}, and explain it.
\end{itemize}

% ---------------------------------------------------------------------
\exerciseset{Options of your own}
\label{c1:set:3}
Create an application named \code{exercise3} to work in for this exercise set.

\exercise{The threshold}{Code}
\label{c1:ex:3.1}
In this exercise you will give the application Kconfig options of its own.

\task{Tasks}
\begin{itemize}
  \item Create a \code{Kconfig} file in \code{exercise3} with a menu named \code{Application}, holding:
\begin{itemize}
  \item \code{APP_HIGH_C}, an \code{int} with the prompt \code{Temperature at which a reading is high, in
    degrees Celsius}, a range from -40 to 125, and a default of 30.
\end{itemize}
  \item End the file so that Zephyr’s own options are still available.
  \item In \code{prj.conf}, set \code{APP_HIGH_C} to 25.
  \item In \code{src/main.c}, go through these readings, printing each with \code{printk()}, and mark those
    at or above \code{CONFIG_APP_HIGH_C} with \code{ (high)}:
\begin{ccode}
static const int readings[] = {-18, 21, 24, 28, 33};
\end{ccode}
\end{itemize}

The output must be:

\begin{console}
-18 C
21 C
24 C
28 C (high)
33 C (high)
\end{console}

\task{Reflection}
\begin{itemize}
  \item Find \code{APP_HIGH_C} in \code{west build -t menuconfig}. What does \code{?} show about it?
  \item Set \code{CONFIG_APP_HIGH_C=200} in \code{prj.conf} and build. What happens, and why is that better
    than the alternative?
  \item Remove the last line of your \code{Kconfig} file and build with \code{-p}. What happens?
\end{itemize}

\exercise{Fahrenheit, when asked for}{Code}
\label{c1:ex:3.2}

\task{Tasks}
\begin{itemize}
  \item Add a \code{bool} option \code{APP_FAHRENHEIT} to the \code{Application} menu, with the prompt
    \code{Print temperatures in degrees Fahrenheit}.
  \item Create \code{src/convert.h} and \code{src/convert.c} with a function:
\begin{ccode}
int celsius_to_fahrenheit(int celsius);
\end{ccode}
which returns the temperature in degrees Fahrenheit, \textbf{rounded to the nearest degree}, using
only integer arithmetic. For example, 37 °C is 98.6 °F, which rounds to 99, and -18 °C is -0.4 °F,
which rounds to 0.
  \item In \code{CMakeLists.txt}, compile \code{src/convert.c} \textbf{only} when
    \code{CONFIG_APP_FAHRENHEIT} is enabled.
  \item In \code{main()}, when the option is enabled, print each reading in Fahrenheit, as \code{<t> F},
    still marked \code{ (high)} by its value in Celsius. Use \code{IS_ENABLED()}, not \code{\#ifdef}.
  \item Create a file \code{fahrenheit.conf} in \code{exercise3} that enables the option.
\end{itemize}

Built with \code{-DEXTRA_CONF_FILE=fahrenheit.conf}, the output must be:

\begin{console}
0 F
70 F
75 F
82 F (high)
91 F (high)
\end{console}

\noindent Built without it, the output must be unchanged from Exercise~\ref{c1:ex:3.1}.

\note[Hint] Work in tenths of a degree: 37 °C is 986 tenths of a degree Fahrenheit. C’s integer
division truncates towards zero, so \code{(tenths + 5) / 10} rounds a positive number of tenths to
the nearest degree, but not a negative one.

\task{Reflection}
\begin{itemize}
  \item Without \code{fahrenheit.conf}, \code{convert.c} is not compiled, yet \code{main.c} still calls
    \code{celsius_to_fahrenheit()}. Why does the program link?
  \item Would it still link with \code{\#ifdef CONFIG_APP_FAHRENHEIT} around the call instead? Which of the
    two would you rather maintain, and why?
  \item Change \code{IS_ENABLED(CONFIG_APP_FAHRENHEIT)} to \code{IS_ENABLED(CONFIG_APP_FARENHEIT)} and build
    with \code{fahrenheit.conf}. What happens? What would have caught it?
\end{itemize}

% ---------------------------------------------------------------------
\exerciseset{Reading the workspace and the build}
\label{c1:set:4}
No code in this set: answer the questions by looking. Use the build of any exercise above.

\exercise[fillaccent2]{The workspace}{Reflection}
\label{c1:ex:4.1}

\task{Tasks}
Answer, using west’s commands rather than by opening files:

\begin{itemize}
  \item Which revision of Zephyr is the workspace using, and where is it checked out?
  \item Which configuration value makes \code{west build} work without \code{-b}?
\end{itemize}

\exercise[fillaccent2]{The build}{Reflection}
\label{c1:ex:4.2}

\task{Tasks}
Answer, by looking in the build directory:

\begin{itemize}
  \item Which log mode (\code{CONFIG_LOG_MODE_...}) did the \code{exercise2} build use? Why that one?
  \item What is the value of \code{CONFIG_BOARD}, and what is the value of \code{CONFIG_BOARD_TARGET}? Why
    are they different?
  \item How many \code{\#define} lines does \code{autoconf.h} hold? Does it contain the options that are
    \code{n}?
  \item Find the definition of \code{k_msleep()} with \code{grep -n -A3 "k_msleep(int32_t ms)"
    deps/zephyr/include/zephyr/kernel.h}. What does it call, and what does \code{static inline} mean for
    whether \code{k_msleep} exists as a function in your program?
\end{itemize}
